% Capítulo 3: Metodología
% Propósito: Detallar el enfoque metodológico, diseño de investigación,
% técnicas e instrumentos de recolección de información y la metodología
% de desarrollo de software seleccionada.

\section{Metodología}
\label{sec:metodologia}

\subsubsection{Enfoque metodológico y tipología del artefacto}
\label{sec:metodo-enfoque}

Para orientar el desarrollo de TravelDLS se adopta como enfoque general la Metodología de Investigación en Ciencia del Diseño, conocida por sus siglas en inglés como DSRM (\textit{Design Science Research Methodology}), propuesta por \textcite{peffers2007dsrm}. Este enfoque resulta pertinente para el proyecto porque la investigación no se limita a describir las dificultades relacionadas con la administración y el seguimiento del transporte terrestre de carga pesada, sino que también busca desarrollar una solución informática que responda a las necesidades identificadas.

La metodología DSRM permite relacionar la identificación de un problema con la definición de los objetivos de la solución, su diseño, construcción y posterior comprobación. Dentro de TravelDLS, esta relación es importante porque permite mantener una conexión entre la problemática que motiva el proyecto y las funcionalidades que se incorporarán a la plataforma. De esta manera, las decisiones técnicas no se consideran actividades aisladas, sino parte de un proceso orientado a responder a las necesidades de gestión, localización, seguimiento y planificación de rutas.

En el ámbito de los sistemas de información, \textcite{hevner2004design} explican que la investigación basada en Ciencia del Diseño puede generar diferentes tipos de artefactos, entre ellos constructos, modelos, métodos e instanciaciones. Estos elementos representan distintas formas de conocimiento que contribuyen a comprender un problema y desarrollar soluciones tecnológicas. Para TravelDLS, el artefacto principal previsto corresponde a una instanciación, representada por la plataforma web que integrará las funciones de administración de empresas transportistas, conductores, vehículos, cargas, órdenes, pagos y seguimiento de los servicios.

La construcción de esta instanciación se apoya en tres elementos complementarios. Los constructos proporcionan los conceptos y el vocabulario necesarios para comprender el dominio de aplicación. Los modelos permiten representar la estructura de la información y el comportamiento esperado del sistema. Por su parte, los métodos comprenden los procedimientos que se utilizarán para organizar la investigación, desarrollar las funcionalidades y resolver problemas específicos, como la planificación de recorridos con múltiples destinos.

Esta organización permite que los fundamentos desarrollados en el marco teórico tengan una aplicación dentro del marco metodológico, evitando repetir definiciones y manteniendo una relación entre los conceptos estudiados, los modelos elaborados y las decisiones que orientarán el desarrollo de TravelDLS.

\subsubsection{Fase 1. Identificación del Problema y Motivación}
\label{sec:metodo-fase1}

La primera fase de la metodología DSRM se orienta a identificar el problema que motiva la investigación y fundamentar la necesidad de desarrollar una solución tecnológica \parencite{peffers2007dsrm}. En TravelDLS, esta etapa parte del análisis de las actividades que intervienen en la gestión del transporte terrestre de carga pesada y de las dificultades que surgen cuando la información necesaria para desarrollar un servicio no se encuentra debidamente relacionada.

El transporte de carga pesada comprende diferentes actividades que deben coordinarse para completar una operación. Además del traslado físico de la mercancía, es necesario administrar solicitudes, vehículos, conductores, cargas, recorridos y comunicaciones entre los participantes del servicio. Estos elementos mantienen una relación directa, debido a que una operación requiere identificar quién contrata el transporte, qué empresa prestará el servicio, cuáles son los recursos asignados y en qué situación se encuentra el traslado.

Cuando esta información se administra mediante registros separados, hojas de cálculo, llamadas o mensajes individuales, resulta más difícil mantener una visión organizada del servicio. Para consultar el avance de una operación es necesario reunir datos provenientes de diferentes medios, lo que dificulta identificar con rapidez su estado y los recursos asociados. Esta situación evidencia la importancia de establecer mecanismos que permitan centralizar la información y mantener una relación entre las diferentes actividades operativas.

En Nicaragua existen investigaciones que han abordado necesidades relacionadas con la administración del transporte mediante herramientas informáticas. \textcite{aragon2018implementacion} desarrollaron un sistema de información para la administración de fletes terrestres, orientado a mejorar la organización de los registros y procesos correspondientes a una empresa de transporte. Este antecedente permite reconocer la importancia que tiene la administración de la información dentro de las operaciones logísticas.

De manera complementaria, \textcite{garcia2022desarrollo} abordaron el seguimiento y la localización de unidades de transporte público en la ciudad de León mediante una aplicación móvil. Aunque este trabajo corresponde a un contexto diferente al transporte de carga pesada, permite reconocer la utilidad de las tecnologías de localización para consultar información sobre el desplazamiento de las unidades.

Estos antecedentes contribuyen a comprender dos necesidades relacionadas con TravelDLS. La primera corresponde a la organización de los registros y procesos administrativos de las empresas transportistas. La segunda se relaciona con la posibilidad de consultar información sobre la ubicación y el desplazamiento de los vehículos durante la prestación de un servicio.

Desde la perspectiva administrativa, una empresa dedicada al transporte de carga pesada necesita mantener relacionados los datos de los recursos que utiliza y las solicitudes que recibe. No resulta suficiente disponer de un registro de conductores y vehículos si esa información no puede relacionarse con las órdenes de transporte que se encuentran pendientes o en ejecución. Por esta razón, la propuesta contempla que las operaciones se organicen mediante una \textbf{\hyperlink{concepto-orden}{orden de transporte}}, que permita vincular los datos del servicio con la empresa, el cliente y los recursos correspondientes.

Asimismo, las solicitudes requieren pasar por diferentes etapas, desde su creación hasta la conclusión del servicio. Esto introduce la necesidad de establecer un \textbf{\hyperlink{concepto-estados}{ciclo de atención}} \textbf{\hyperlink{concepto-estados}{y estados de una orden}}, mediante el cual se identifique cuándo una solicitud se encuentra pendiente, cuándo ha recibido un precio y cuándo ha sido aceptada o rechazada por el cliente. La definición de estos estados permitirá representar el avance administrativo y operativo de una solicitud, manteniendo una relación entre las acciones realizadas y la situación actual del servicio.

Otra dificultad relacionada con la administración del transporte corresponde a la \textbf{\hyperlink{concepto-asignacion}{asignación de recursos}}. Cada solicitud requiere identificar los medios que participarán en su atención, principalmente el conductor y el vehículo. Cuando esta relación no se mantiene dentro de un mismo registro, la consulta de la operación depende de información adicional que puede encontrarse en medios diferentes. TravelDLS busca organizar estos datos para que la empresa pueda identificar los recursos asociados a las órdenes que administra.

Desde la perspectiva del cliente, una de las principales necesidades consiste en disponer de información sobre el servicio contratado. Después de aceptar una solicitud de transporte, resulta importante conocer si la carga ha comenzado a trasladarse, cuál es el estado del servicio y dónde se encuentra la unidad encargada del recorrido.

Cuando esta información no puede consultarse de manera directa, el cliente depende de la comunicación individual con la empresa o el conductor para conocer el avance de su carga. Esto puede generar incertidumbre sobre la situación del traslado, especialmente cuando no existen mecanismos que permitan consultar la ubicación reportada y los cambios de estado de la operación.

La incorporación de funciones de \textbf{\hyperlink{concepto-trazabilidad}{trazabilidad}} y \textbf{\hyperlink{concepto-seguimiento-tiempo-real}{seguimiento en tiempo real}} responde a esta necesidad. Mediante estas funcionalidades se pretende mantener vinculados los registros de ubicación y los estados de la operación con la orden correspondiente, permitiendo que los usuarios autorizados consulten información relacionada con el desarrollo del servicio.

La \textbf{\hyperlink{concepto-geolocalizacion}{geolocalización}} constituye un elemento necesario para representar la posición de la unidad durante el traslado. Sin embargo, conocer únicamente la ubicación más reciente no permite reconstruir los desplazamientos realizados anteriormente. Por esta razón, también se contempla el \textbf{\hyperlink{concepto-historial}{historial de seguimiento}}, que permitirá conservar registros sucesivos asociados a una misma operación.

Esta diferencia resulta importante porque el seguimiento de un servicio no depende solamente de mostrar un punto sobre un mapa. También requiere relacionar la ubicación reportada con el conductor, la orden correspondiente y el estado del traslado. De esta manera, la información consultada podrá interpretarse dentro del contexto de la operación que se está desarrollando.

No obstante, el seguimiento estará condicionado por la disponibilidad de conexión a Internet, los permisos de localización y las características de los dispositivos utilizados. Estas condiciones deben considerarse al establecer las funcionalidades y limitaciones del proyecto, evitando afirmar que la plataforma garantizará una actualización continua de ubicaciones bajo cualquier circunstancia.

Además de la administración y el seguimiento, otra necesidad identificada corresponde a la organización de recorridos con múltiples destinos. Una empresa transportista puede requerir visitar diferentes puntos durante un mismo servicio, por lo que resulta pertinente analizar el orden en que se realizará el desplazamiento.

La \textbf{\hyperlink{concepto-planificacion-rutas}{planificación de rutas}} organiza los destinos. La \mbox{\textbf{\hyperlink{concepto-optimizacion-rutas}{optimización de rutas}}} permite comparar recorridos según criterios como la distancia o el tiempo. Esto resulta relevante debido a que el orden de visita influye en la distancia total que puede representar una operación.


En el contexto nacional, \textcite{aragon2018implementacion}
abordaron la administración de fletes terrestres mediante
un sistema de información, mientras que
\textcite{garcia2022desarrollo} estudiaron el seguimiento
y localización de unidades de transporte. Aunque estas
investigaciones no se centran directamente en la
optimización de rutas de carga pesada, permiten
reconocer la aplicación de herramientas tecnológicas
en el sector del transporte nicaragüense.

Por otra parte, los antecedentes regionales de
\textcite{aguilar2022optimizacion} y
\textcite{hernandez2024evaluacion} abordan la
planificación y optimización de rutas mediante
diferentes alternativas. Estas investigaciones
constituyen una referencia para el componente de
planificación de recorridos que se contempla
desarrollar en TravelDLS.


Estas investigaciones permiten fundamentar la pertinencia de estudiar alternativas de optimización, aunque los resultados obtenidos corresponden a los escenarios particulares de cada trabajo y no pueden considerarse mejoras garantizadas para TravelDLS.

Por esta razón, el proyecto contempla analizar el \textbf{\hyperlink{concepto-tsp}{Problema del Agente Viajero}} y comparar procedimientos como \textbf{\hyperlink{concepto-vecino-mas-cercano}{Vecino Más Cercano}} y \textbf{\hyperlink{concepto-dos-opt}{2-opt}} con las soluciones generadas mediante \textbf{\hyperlink{concepto-ortools}{Google OR-Tools}}. El propósito de esta comparación será conocer el comportamiento de las alternativas consideradas y establecer una base técnica para el módulo de planificación de rutas del prototipo.

La problemática también se relaciona con la presencia tecnológica de las empresas transportistas. TravelDLS contempla reunir diferentes organizaciones dentro de una misma plataforma, permitiendo que puedan presentar sus servicios a clientes interesados en contratar transporte de carga pesada. Esta característica busca conectar la administración interna con las actividades de búsqueda y contratación, ofreciendo un espacio común para empresas y clientes.

La posibilidad de mejorar la visibilidad de las empresas constituye un beneficio esperado de la propuesta. Sin embargo, no debe presentarse como un resultado demostrado, debido a que todavía sería necesario evaluar el uso de la plataforma y su aceptación dentro del entorno operativo correspondiente.

Como respuesta a las necesidades identificadas, TravelDLS propone desarrollar una plataforma bajo el modelo \textbf{\hyperlink{concepto-saas}{Software como Servicio}}, destinada a integrar información administrativa, operativa y de seguimiento dentro de un mismo entorno. La \textbf{\hyperlink{concepto-multitenant}{arquitectura multitenant}} y la \textbf{\hyperlink{concepto-separacion}{separación de información operativa}} permitirán orientar el diseño para que diferentes empresas puedan utilizar la plataforma manteniendo diferenciados sus registros y las funciones correspondientes a cada perfil.

De esta manera, la motivación principal del proyecto consiste en ofrecer una alternativa tecnológica que permita relacionar las operaciones de transporte, mantener información accesible para los participantes autorizados y apoyar la planificación de recorridos. TravelDLS no busca sustituir las decisiones administrativas de las empresas, sino proporcionar una herramienta que contribuya a organizar los procesos y facilitar la consulta de información durante la prestación del servicio.

La identificación del problema permite establecer una base para definir las características que deberá cumplir la solución. No obstante, las dificultades descritas a partir del análisis documental y los antecedentes deberán diferenciarse de aquellas que puedan comprobarse posteriormente mediante información obtenida directamente del entorno operativo.

En consecuencia, esta primera fase delimita las necesidades que motivan TravelDLS y proporciona los fundamentos para formular los objetivos de la solución, sin considerar como demostrados los beneficios que todavía requieren comprobación.


% =========================================================
% FASE 2
% =========================================================

\subsubsection{Fase 2. Definición de Objetivos de la Solución (Meta-requisitos)}
\label{sec:metodo-fase2}

La segunda fase de la metodología DSRM consiste en establecer las características que deberá cumplir la solución propuesta para responder a las necesidades identificadas durante la etapa anterior \parencite{peffers2007dsrm}. En TravelDLS, esta fase permite relacionar la problemática con los objetivos del proyecto y formular los meta-requisitos que orientarán el diseño de la plataforma.

Los meta-requisitos representan las capacidades generales que se espera que el sistema proporcione. Su formulación permite describir qué necesita resolver TravelDLS antes de establecer cómo serán implementadas técnicamente sus funcionalidades. Esta diferencia es importante porque permite separar las necesidades del sistema de las herramientas, estructuras y procedimientos que se utilizarán posteriormente para atenderlas.

A partir del problema identificado, uno de los principales meta-requisitos corresponde a la centralización de la información administrativa y operativa. La plataforma deberá permitir registrar y consultar los datos relacionados con empresas, clientes, conductores, vehículos y órdenes de transporte, manteniendo las asociaciones necesarias entre estos elementos.

Otro meta-requisito se relaciona con la administración de las solicitudes de transporte. TravelDLS deberá permitir que el cliente registre una solicitud, que la empresa establezca el precio correspondiente y que el cliente pueda aceptar o rechazar las condiciones del servicio. Asimismo, será necesario mantener un estado que permita identificar la situación actual de cada orden.

En relación con el seguimiento, la solución deberá permitir vincular los registros de ubicación y estado con las órdenes correspondientes, de manera que los usuarios autorizados puedan consultar información sobre el desplazamiento de sus cargas. Esta capacidad deberá considerar las limitaciones de conectividad y disponibilidad de los dispositivos utilizados.

La plataforma también deberá contemplar la asignación de recursos de transporte, permitiendo relacionar las operaciones con los conductores y vehículos correspondientes. Esta característica contribuirá a mantener organizada la información necesaria para identificar los recursos involucrados en cada servicio.

En cuanto a la planificación de rutas, TravelDLS deberá proporcionar un mecanismo que permita analizar recorridos con múltiples destinos. Para este componente se contempla la comparación de Vecino Más Cercano y 2-opt con Google OR-Tools, utilizando criterios equivalentes para estudiar sus resultados. La integración del módulo de optimización se orientará a apoyar la planificación, sin sustituir las decisiones que correspondan a la empresa transportista.

Asimismo, el sistema deberá mantener la separación de la información perteneciente a diferentes empresas transportistas, estableciendo las restricciones necesarias para que los usuarios accedan únicamente a los datos relacionados con sus funciones y servicios autorizados.

Finalmente, el proceso de contratación deberá mostrar el precio final del servicio antes de que el cliente lo acepte, incorporando el porcentaje destinado al mantenimiento de la plataforma. Esta condición responde al modelo de funcionamiento económico definido para TravelDLS y busca que el cliente conozca el monto correspondiente antes de confirmar la contratación.

Estos meta-requisitos constituyen una propuesta inicial que relaciona las necesidades identificadas con las capacidades esperadas del sistema. Su organización permitirá orientar las decisiones de diseño y establecer posteriormente los criterios necesarios para comprobar el funcionamiento de cada componente.

La siguiente matriz presenta esta relación inicial, sin asignar porcentajes de mejora o resultados cuantitativos que todavía no hayan sido determinados mediante procedimientos de evaluación.



% =========================================================
% MATRIZ DE META-REQUISITOS DE TRAVELDLS
% =========================================================

\begin{table}[H]
	\centering
	\caption{Matriz inicial de meta-requisitos de TravelDLS}
	\label{tab:meta-requisitos-traveldls}
	
	\small
	\setlength{\tabcolsep}{4.5pt}
	\renewcommand{\arraystretch}{1.45}
	
	\begin{tabularx}{\textwidth}{
			@{}
			>{\RaggedRight\arraybackslash}p{1.35cm}
			>{\RaggedRight\arraybackslash}p{3.05cm}
			>{\RaggedRight\arraybackslash}X
			>{\RaggedRight\arraybackslash}p{2.85cm}
			@{}
		}
		
		\toprule
		
		\textbf{Código} &
		\textbf{Necesidad identificada} &
		\textbf{Meta-requisito previsto} &
		\textbf{Relación con el diseño} \\
		
		\midrule
		
		MR-01 &
		Organización de información administrativa y operativa. &
		Registrar y consultar clientes, empresas, vehículos, conductores y órdenes desde una plataforma compartida. &
		Figuras~\ref{fig:diagrama-clases} y~\ref{fig:diagrama-er}. \\
		
		\midrule
		
		MR-02 &
		Administración de solicitudes de transporte. &
		Permitir la creación, cotización, aceptación, rechazo y consulta de estados de una orden. &
		Figuras~\ref{fig:diagrama-estados} y~\ref{fig:diagrama-secuencia}. \\
		
		\midrule
		
		MR-03 &
		Consulta de información sobre las cargas. &
		Asociar registros de ubicación y estado a las órdenes, permitiendo su consulta por usuarios autorizados. &
		Figuras~\ref{fig:diagrama-clases},~\ref{fig:diagrama-er} y~\ref{fig:diagrama-secuencia}. \\
		
		\midrule
		
		MR-04 &
		Coordinación de recursos de transporte. &
		Relacionar conductores y vehículos con las operaciones según las reglas establecidas. &
		Figuras~\ref{fig:diagrama-clases} y~\ref{fig:diagrama-er}. \\
		
		\midrule
		
		MR-05 &
		Planificación de recorridos con múltiples destinos. &
		Comparar métodos de optimización con entradas equivalentes y desarrollar el módulo previsto para el prototipo. &
		\textbf{\hyperref[sec:objetivos-especificos]{Objetivos del proyecto}} y módulo de optimización. \\
		
		\midrule
		
		MR-06 &
		Claridad en el precio del servicio. &
		Mostrar el monto final de la contratación, incluyendo el porcentaje de mantenimiento, antes de su aceptación. &
		Figuras~\ref{fig:diagrama-estados} y~\ref{fig:diagrama-secuencia}. \\
		
		\midrule
		
		MR-07 &
		Uso de la plataforma por diferentes empresas. &
		Mantener diferenciados los registros operativos de cada empresa y establecer controles de acceso según los perfiles. &
		Figuras~\ref{fig:diagrama-clases} y~\ref{fig:diagrama-er}. \\
		
		\bottomrule
		
	\end{tabularx}
	
	\vspace{0.2cm}
	
	\begin{minipage}{\textwidth}
		\footnotesize
		\textit{Nota.} Los meta-requisitos representan capacidades
		previstas para TravelDLS. Su cumplimiento deberá comprobarse
		posteriormente mediante criterios de aceptación y las
		evidencias correspondientes. La matriz no representa
		resultados de una evaluación finalizada.
	\end{minipage}
	
\end{table}

\subsubsection{Fase 3. Diseño y Desarrollo del Artefacto (Meta-diseños)}
\label{sec:metodo-fase3}

La tercera fase de la metodología DSRM comprende el diseño y la construcción de la solución tecnológica a partir de los objetivos establecidos previamente \parencite{peffers2007dsrm}. En TravelDLS, esta etapa permite relacionar los meta-requisitos con las decisiones de arquitectura, los modelos de información y los componentes que conformarán la plataforma web.

Los meta-diseños representan las decisiones técnicas destinadas a satisfacer las capacidades esperadas del sistema. En este proyecto, dichas decisiones se encuentran relacionadas con la organización de los módulos, el almacenamiento de información, las interacciones entre los usuarios y la definición de los procedimientos que permitirán desarrollar las operaciones de transporte.

Como parte del diseño actual de TravelDLS se dispone de un diagrama de clases y de un diagrama entidad-relación. El primero representa las principales entidades y servicios del sistema, mientras que el segundo permite comprender la estructura prevista para el almacenamiento de datos. Ambos modelos constituyen una base para establecer las relaciones necesarias entre usuarios, empresas, clientes, conductores, vehículos, órdenes, detalles de carga, pagos y registros de seguimiento.

El modelo de clases también permite reconocer los servicios destinados a administrar operaciones específicas, como la autenticación de usuarios, la gestión de órdenes y el procesamiento de pagos. Estas responsabilidades contribuyen a establecer una separación entre los componentes que representan los datos y aquellos que participan en la ejecución de determinadas funciones.

El diagrama de clases presentado en la Figura~\ref{fig:diagrama-clases} organiza los componentes principales del sistema alrededor de usuarios, roles, empresas, clientes, conductores, vehículos, órdenes, detalles de carga, pagos y seguimiento. Esta representación permite observar cómo los servicios de autenticación, órdenes y pagos se relacionan con las entidades del dominio, de manera que las operaciones previstas mantengan una separación entre responsabilidades funcionales y estructuras de información.

El diagrama entidad-relación de la Figura~\ref{fig:diagrama-er} complementa el modelo anterior al presentar las tablas, claves y asociaciones previstas para la persistencia de datos. Su incorporación permite relacionar el diseño conceptual con la estructura de almacenamiento, especialmente en los elementos asociados a empresas, clientes, conductores, vehículos, órdenes, detalles de carga, pagos y registros de ubicación.

En relación con el comportamiento del sistema, el diagrama de estados permite representar las situaciones que puede experimentar una orden de transporte. Su utilización contribuye a establecer las transiciones necesarias durante la creación, cotización, aceptación, ejecución y finalización de una solicitud.

Como se muestra en la Figura~\ref{fig:diagrama-estados}, la orden inicia en estado pendiente y avanza según las acciones realizadas por la empresa y el cliente. La asignación de precio, la aceptación o rechazo, la posible anulación, el pago, el inicio del viaje y la entrega establecen reglas que ayudan a mantener coherencia entre el avance administrativo de la solicitud y su ejecución operativa.

Por su parte, el diagrama de secuencia describe la interacción prevista entre el cliente, la empresa transportista, el conductor, el frontend, el backend y la base de datos. Este modelo permite observar cómo se intercambiará información durante los procesos de creación de órdenes, asignación de precios, registro de pagos, asignación de conductores y seguimiento de los servicios.

La Figura~\ref{fig:diagrama-secuencia} muestra el flujo principal de una orden desde su creación hasta el seguimiento del viaje. En este proceso, el cliente registra la solicitud, la empresa asigna el precio, el cliente acepta la cotización, se registra el pago, se asigna el conductor y posteriormente se consultan los registros de seguimiento. Esta secuencia permite vincular las operaciones de interfaz, backend y base de datos con los estados definidos para la orden.

La relación entre ambos diagramas es importante porque las operaciones representadas en las secuencias deben mantener coherencia con los estados definidos para las órdenes. Por esta razón, durante la consolidación del diseño será necesario revisar las reglas relacionadas con el momento de procesamiento del pago, la asignación del conductor y el inicio del traslado, evitando que un mismo proceso tenga interpretaciones diferentes entre los modelos.

En cuanto a la organización tecnológica, el diseño documentado contempla una plataforma web cuyo frontend utiliza Next.js y cuyo backend se desarrolla con AdonisJS. Estos componentes permiten separar la interfaz utilizada por los usuarios de los procedimientos encargados de procesar solicitudes, aplicar las reglas de negocio y administrar la información almacenada en la base de datos.

\paragraph{Diseño propuesto para la optimización de rutas}
El módulo de optimización de TravelDLS aún no está implementado. Como parte del meta-diseño se propone incorporar Google OR-Tools en una etapa posterior, con entradas y restricciones documentadas. Se distinguen dos ejercicios: un trayecto directo entre una recogida A y una entrega B, utilizado solo para ilustrar la elección de un camino en una red, y un recorrido que parte de A después de recoger la carga y visita varios destinos. Este segundo caso corresponde a la planificación con múltiples destinos prevista en los meta-requisitos y en el plan de ejecución.

Para el trayecto directo se propone representar la red vial mediante un grafo dirigido $G=(V,E)$, con un costo no negativo $c_{ij}$ para cada arco disponible. Una posible formulación en \texttt{SimpleMinCostFlow} asigna una unidad de suministro a A, una unidad de demanda a B, cero a los nodos intermedios y capacidad unitaria a cada arco \parencite{googleMinCostFlow}. La conservación del flujo se expresa como $\sum_j x_{vj}-\sum_i x_{iv}=b_v$, donde $b_A=1$, $b_B=-1$ y $b_v=0$ para los demás nodos; el objetivo es minimizar $\sum_{(i,j)\in E}c_{ij}x_{ij}$. Con costos enteros positivos, la unidad de flujo permite identificar un camino de costo mínimo. Las conexiones, sentidos de circulación y restricciones para vehículos de carga deberán obtenerse de datos viales adecuados; OR-Tools no los aporta por sí solo. Este ejercicio de camino mínimo no sustituye la comparación de métodos para múltiples destinos.

Para varias entregas se propone \texttt{RoutingModel}, con A como inicio fijo, una matriz común de distancias enteras y la obligación de visitar todos los destinos \parencite{googleTSP}. El ejercicio académico supondrá el retorno a A, conforme al TSP descrito en el marco teórico; antes de aplicar el modelo a un viaje real habrá que definir si el vehículo efectivamente regresa al origen y considerar las restricciones operativas pertinentes. Si se añaden recogidas y entregas emparejadas, deberán imponerse además las condiciones de mismo vehículo y precedencia, no solo registrar los pares mediante \texttt{AddPickupAndDelivery} \parencite{googlePickupDelivery}.

La comparación prevista para múltiples destinos utilizará la misma matriz, origen, destinos y condición de retorno para Vecino Más Cercano, su mejora mediante 2-opt y OR-Tools. Se registrarán el orden de visita, la distancia total y el tiempo de cálculo bajo condiciones de ejecución documentadas. Los ejercicios hipotéticos del apartado de resultados permiten verificar operaciones matemáticas, pero no constituyen pruebas de funcionamiento ni mediciones de rendimiento del módulo futuro. Tampoco se presupone que el solucionador de rutas encuentre siempre el óptimo global \parencite{googleTSP}.

\subsection{Metodología de Desarrollo}
\label{sec:metodologia-desarrollo}

La metodología de desarrollo de TravelDLS se integra en la fase 3 de la metodología de investigación en ciencia del diseño (Design Science Research Methodology, DSRM), denominada Diseño y Desarrollo del Artefacto. En esta fase, las decisiones de diseño y arquitectura derivadas de los meta-requisitos se utilizan como base para construir la plataforma web. Para organizar este proceso, se propone combinar Scrum y Extreme Programming (XP). Scrum se utilizará para planificar y distribuir el trabajo de forma incremental, mientras que XP orientará las prácticas de ingeniería hacia la calidad técnica del software \parencite{beck2000xp, schwaber2020scrum}.

Esta combinación se enmarca en el enfoque de investigación adoptado para TravelDLS. \textcite{peffers2007dsrm} describen la metodología DSRM como un proceso que comprende la identificación del problema, la definición de los objetivos de la solución, el diseño y desarrollo del artefacto, su demostración, su evaluación y la comunicación de los resultados. En consecuencia, el nivel táctico se concentra en la construcción del sistema a partir de los meta-requisitos establecidos, sin confundir la implementación de las funcionalidades con los resultados que deberán obtenerse durante su evaluación.

La metodología de desarrollo describe cómo se organizará la construcción del artefacto TravelDLS a partir de los siete meta-requisitos definidos en la Tabla~\ref{tab:meta-requisitos-traveldls}. Esta actividad se relaciona con la fase de diseño y desarrollo del artefacto dentro de la metodología de investigación en ciencia del diseño (DSRM).\textcite{peffers2007dsrm} presentan la DSRM como un proceso que vincula la identificación del problema, los objetivos de la solución, el diseño y desarrollo, la demostración, la evaluación y la comunicación de resultados. Por ello, las matrices de esta sección conectan las necesidades del sistema con las decisiones de diseño, las historias de usuario y las formas previstas de comprobar su cumplimiento.

Para organizar el trabajo se propone Scrum como marco de gestión y se consideran prácticas técnicas de Extreme Programming (XP) para apoyar la calidad del código. \textcite{bahit2012scrum} explica cómo pueden combinarse ambos enfoques, utilizando Scrum para organizar el trabajo en iteraciones y XP para reforzar las prácticas de desarrollo. La aplicación de estas prácticas en TravelDLS deberá verificarse durante la implementación y mediante las pruebas correspondientes \parencite{beck2000xp, bahit2012scrum, schwaber2020scrum}.

\subsubsection{Gestión del desarrollo mediante Scrum}

Scrum es un marco de trabajo que permite abordar problemas complejos mediante un proceso iterativo e incremental. Se fundamenta en el empirismo y el pensamiento Lean, y establece la transparencia, la inspección y la adaptación como pilares para gestionar el trabajo. Estos principios permiten revisar el avance, identificar desviaciones y ajustar las actividades según la información obtenida durante el desarrollo \parencite{schwaber2020scrum}.

Para TravelDLS, se propone utilizar Scrum para organizar las actividades de desarrollo mediante un Product Backlog, historias de usuario y Sprints. El Product Backlog reunirá las funcionalidades pendientes del sistema y se actualizará conforme se definan nuevas necesidades o se precise el alcance de las ya identificadas. Las historias de usuario describirán las funcionalidades desde la perspectiva de los perfiles que utilizarán la plataforma, manteniendo su relación con los meta-requisitos establecidos. El Product Backlog reúne las historias de usuario priorizadas según su valor funcional y las dependencias entre tareas. En cada sprint, el equipo selecciona el trabajo que abordará y revisa el incremento obtenido; la retrospectiva permite identificar ajustes para la siguiente iteración (Bahit, 2012; Schwaber \& Sutherland, 2020). La asignación de historias a los sprints es provisional y debe concordar con el \mbox{cronograma vigente}.

El trabajo se organizará en Sprints de dos a tres semanas, según la planificación propuesta en el documento metodológico. Al inicio de cada Sprint se establecerá un objetivo y se seleccionarán los elementos del Product Backlog que se atenderán, considerando su prioridad, complejidad y dependencias técnicas. Durante la ejecución se desarrollarán las funcionalidades previstas y se realizarán las comprobaciones correspondientes. Al finalizar cada Sprint, se revisará el incremento obtenido y se identificarán ajustes para las iteraciones siguientes, mediante actividades de revisión y retrospectiva (Schwaber \& Sutherland, 2020).

La planificación tendrá en cuenta la arquitectura tecnológica definida para TravelDLS: Next.js para el frontend, AdonisJS para el backend y PostgreSQL para la persistencia de datos. El backend se organizará mediante el patrón MVC y una capa de servicios para separar las responsabilidades de los controladores y la lógica de negocio. La distribución de actividades considerará las dependencias entre los componentes, de modo que la implementación de las interfaces se integre con los servicios y las estructuras de datos correspondientes.

\subsubsection{Prácticas de ingeniería mediante Extreme Programming (XP)}

Extreme Programming (XP) propone prácticas de ingeniería que buscan mejorar la calidad del software y facilitar su adaptación durante el desarrollo. \textcite{beck2000xp} describe este enfoque a partir de prácticas que promueven la retroalimentación frecuente, la simplicidad del diseño y la colaboración entre los integrantes del equipo. En TravelDLS se propone incorporar cuatro prácticas: desarrollo guiado por pruebas, refactorización continua, programación en pares y propiedad colectiva del código, e integración continua.

\begin{itemize}
	\item \textbf{Desarrollo guiado por pruebas (TDD):} El desarrollo guiado por pruebas (\textit{Test-Driven Development}, TDD) consiste en escribir una prueba automatizada antes del código funcional que permitirá satisfacerla. A partir del resultado de la prueba, se implementa la funcionalidad y se realizan los ajustes necesarios hasta que se cumpla el comportamiento esperado \parencite{beck2000xp}. En TravelDLS se propone aplicar esta práctica principalmente a las reglas de negocio del backend, como las validaciones de usuarios, la asignación de vehículos y conductores, la gestión de órdenes y el procesamiento de información relacionada con los recorridos. Las pruebas permitirán comprobar el comportamiento esperado de cada componente y detectar errores durante su implementación.
\end{itemize}

\begin{itemize}
	\item \textbf{Refactorización continua:} La refactorización consiste en mejorar la estructura interna del código sin modificar su comportamiento observable \parencite{beck2000xp}. En TravelDLS se propone revisar y reorganizar el código conforme se incorporen nuevas funcionalidades, especialmente cuando existan instrucciones repetidas, responsabilidades mezcladas o procedimientos que puedan simplificarse. Esta práctica se aplicará a los componentes del frontend desarrollado con Next.js y a los módulos del backend construido con AdonisJS, procurando conservar la separación entre la interfaz, los controladores, la capa de servicios y el acceso a los datos.
	
	\item \textbf{Programación en pares y propiedad colectiva del código:} La programación en pares consiste en que dos integrantes colaboren en una misma tarea de programación, mientras que la propiedad colectiva permite que los integrantes del equipo puedan revisar y modificar diferentes partes del código, en lugar de limitar cada componente a una sola persona \parencite{beck2000xp}. Para TravelDLS se propone utilizar estas prácticas en funcionalidades que involucren varios componentes del sistema, como la autenticación, la gestión de órdenes y la optimización de recorridos. La colaboración facilitará el intercambio de conocimientos, la revisión de las soluciones implementadas y el mantenimiento del código. Su aplicación se ajustará a la disponibilidad de los integrantes y a la distribución de las tareas de cada sprint.
\end{itemize}

\begin{itemize}
	\item \textbf{Integración continua:} \textcite{bahit2012scrum} explica que el objetivo de la integración continua es incorporar las nuevas funcionalidades a las existentes, procurando que los cambios realizados no afecten el funcionamiento de las funcionalidades anteriores. En TravelDLS, se propone aplicar esta práctica mediante la integración progresiva de los cambios en el código y la ejecución de pruebas que permitan comprobar que las nuevas funcionalidades se incorporen correctamente sin alterar las ya desarrolladas. Esto contribuirá a detectar errores durante el desarrollo y a mantener el funcionamiento conjunto de los componentes del sistema.
\end{itemize}

\subsubsection{Matriz de meta-diseños}

Los meta-diseños resumen las decisiones técnicas previstas para responder a los meta-requisitos; su cumplimiento se comprobará durante el desarrollo.


\begingroup
	\small
	\setlength{\tabcolsep}{4.5pt}
	\renewcommand{\arraystretch}{1.45}
	\begin{longtable}{@{}>{\RaggedRight\arraybackslash}p{1.5cm}>{\RaggedRight\arraybackslash}p{1.8cm}>{\RaggedRight\arraybackslash}p{4.1cm}>{\RaggedRight\arraybackslash}p{2.3cm}>{\RaggedRight\arraybackslash}p{5.3cm}@{}}
		\caption{Matriz de meta-diseños asociados a los meta-requisitos de TravelDLS}\label{tab:meta_disenos}\\
		\toprule
		\textbf{Código} & \textbf{MR asociado} & \textbf{Decisión de diseño} & \textbf{Historias} & \textbf{Verificación} \\
		\midrule
		\endfirsthead
		\caption[]{Matriz de meta-diseños asociados a los meta-requisitos de TravelDLS (continuación)}\\
		\toprule
		\textbf{Código} & \textbf{MR asociado} & \textbf{Decisión de diseño} & \textbf{Historias} & \textbf{Verificación} \\
		\midrule
		\endhead
		\endfoot
		\bottomrule
		\endlastfoot
		MD-01 & MR-01 & Modelo relacional en PostgreSQL para usuarios, roles, empresas, clientes, conductores, vehículos, órdenes y registros asociados; operaciones CRUD mediante AdonisJS, organizadas con MVC y capa de servicios. & HU-01, HU-05, HU-14, HU-29 & Integridad de relaciones; persistencia de altas y cambios; consulta paginada; errores de validación. \\ \midrule
		MD-02 & MR-02 & Flujo de órdenes con estados definidos y operaciones para crear solicitud, cotizar, aceptar o rechazar, consultar y modificar según el estado permitido. & HU-16, HU-20 & Transiciones válidas; cotización visible antes de decidir; rechazo de operaciones incompatibles. \\ \midrule
		MD-03 & MR-03 & Registros de ubicación y estado vinculados a una orden; consulta restringida a participantes autorizados; mapa y ubicación real separados de la simulación del recorrido. & HU-22, HU-24 & Última ubicación asociada al viaje correcto; permisos aplicados; simulación sin modificar GPS ni estado real. \\ \midrule
		MD-04 & MR-04 & Relaciones entre empresas, conductores, camiones y pedidos, con reglas de disponibilidad, asignación, estado operativo y reporte de mantenimiento. & HU-11, HU-15, HU-21, HU-22 & No permitir asignaciones incompatibles; conductor consulta solo su vehículo/viajes; cambios de disponibilidad coherentes. \\ \midrule
		MD-05 & MR-05 & Módulo de recorridos con origen y destinos, matriz de costo, integración de Google OR-Tools y comparación con Vecino Más Cercano y 2-opt; visualización mediante Mapbox. & HU-30, HU-33 & Mismas entradas y condiciones en la comparación; costo y tiempo de cálculo registrados; no afirmar óptimo garantizado. \\ \midrule
		MD-06 & MR-06 & Presentación del precio final antes de la aceptación, incluyendo el porcentaje de mantenimiento previsto en el modelo económico; flujo de pago del prototipo identificado explícitamente como simulado. & HU-16, HU-19, HU-20, HU-25, HU-28 & Importe final visible antes de aceptar; reglas de estado y total consistentes; ninguna operación se presenta como cobro bancario real. \\ \midrule
		MD-07 & MR-07 & Autenticación, sesión y autorización por rol, con comprobación de pertenencia de los registros a la empresa, cliente, conductor o pedido correspondiente. & HU-03, HU-04, HU-06, HU-07, HU-11, HU-14, HU-16, HU-18, HU-21, HU-23, HU-26, HU-29, HU-30, HU-33 & Denegar acceso directo o por API a datos ajenos; mantener separada la información por empresa y perfil. \\ \midrule
	\end{longtable}
\endgroup

\textit{Nota.} Elaboración propia a partir de los meta-requisitos, los diagramas de diseño y las historias de usuario de TravelDLS.

\subsubsection{Matriz de trazabilidad de historias de usuario}

La Tabla~\ref{tab:trazabilidad_tactica} muestra la contribución de cada historia de usuario a los meta-requisitos de la Tabla~\ref{tab:meta-requisitos-traveldls}.


\begingroup
	\small
	\setlength{\tabcolsep}{4.5pt}
	\renewcommand{\arraystretch}{1.45}
	\begin{longtable}{@{}>{\RaggedRight\arraybackslash}p{1.7cm}>{\RaggedRight\arraybackslash}p{3.9cm}>{\RaggedRight\arraybackslash}p{2.8cm}>{\RaggedRight\arraybackslash}p{7cm}@{}}
		\caption{Matriz de trazabilidad entre historias de usuario y meta-requisitos}\label{tab:trazabilidad_tactica}\\
		\toprule
		\textbf{Historia} & \textbf{Resumen} & \textbf{Meta-requisitos} & \textbf{Área funcional} \\
		\midrule
		\endfirsthead
		\caption[]{Matriz de trazabilidad entre historias de usuario y meta-requisitos (continuación)}\\
		\toprule
		\textbf{Historia} & \textbf{Resumen} & \textbf{Meta-requisitos} & \textbf{Área funcional} \\
		\midrule
		\endhead
		\endfoot
		\bottomrule
		\endlastfoot
		HU-01 & Cimientos de base de datos y backend & MR-01 & Preparación de infraestructura y persistencia \\ \midrule
		HU-02 & Registrarse como cliente & MR-01, MR-07 & Registro de clientes y asignación del rol cliente \\ \midrule
		HU-03 & Iniciar, conservar y cerrar sesión & MR-01, MR-07 & Autenticación y ciclo de sesión \\ \midrule
		HU-04 & Acceder a navegación y operaciones por rol & MR-07 & Navegación y autorización por perfil \\ \midrule
		HU-05 & Registrar usuarios con rol desde administración & MR-01, MR-07 & Administración de usuarios \\ \midrule
		HU-06 & Administrar empresas & MR-01, MR-07 & Gestión de empresas \\ \midrule
		HU-07 & Consultar y dar de baja clientes desde administración & MR-01, MR-07 & Consulta y baja de clientes \\ \midrule
		HU-08 & Completar y actualizar perfil de cliente & MR-01, MR-07 & Perfil de cliente \\ \midrule
		HU-09 & Actualizar perfil y tarifa de empresa & MR-01, MR-06, MR-07 & Perfil y tarifa de empresa \\ \midrule
		HU-10 & Completar y actualizar perfil de conductor & MR-01, MR-07 & Perfil de conductor \\ \midrule
		HU-11 & Registrar, vincular y administrar conductores & MR-01, MR-04, MR-07 & Gestión y vinculación de conductores \\ \midrule
		HU-12 & Administrar categorías de vehículos & MR-01, MR-04 & Categorías de vehículos \\ \midrule
		HU-13 & Administrar camiones y asignarlos a conductores & MR-01, MR-04, MR-07 & Gestión de camiones y asignación de conductores \\ \midrule
		HU-14 & Consultar camión asignado & MR-01, MR-04, MR-07 & Consulta del camión asignado \\ \midrule
		HU-15 & Reportar mantenimiento del camión & MR-04 & Reporte de mantenimiento \\ \midrule
		HU-16 & Solicitar transporte con empresa, carga y direcciones & MR-02, MR-07 & Solicitud de transporte \\ \midrule
		HU-17 & Consultar y gestionar pedidos del cliente & MR-02, MR-07 & Consulta y cambios permitidos de pedidos \\ \midrule
		HU-18 & Consultar solicitudes recibidas por empresa & MR-02, MR-07 & Consulta de solicitudes por empresa \\ \midrule
		HU-19 & Establecer el precio de un pedido & MR-02, MR-07 & Cotización del servicio \\ \midrule
		HU-20 & Responder a la cotización de transporte & MR-02, MR-07 & Aceptación o rechazo de cotización \\ \midrule
		HU-21 & Asignar conductor a un pedido & MR-04, MR-07 & Asignación de conductor al pedido \\ \midrule
		HU-22 & Consultar, iniciar y finalizar viajes & MR-03, MR-04, MR-07 & Gestión de viajes \\ \midrule
		HU-23 & Compartir y consultar ubicación del viaje & MR-03, MR-07 & Ubicación real y consulta del seguimiento \\ \midrule
		HU-24 & Visualizar simulación del recorrido & MR-03, MR-07 & Simulación visual separada del seguimiento real \\ \midrule
		HU-25 & Realizar pago simulado por tarjeta o transferencia & MR-02, MR-07 & Pago simulado de una cotización \\ \midrule
		HU-26 & Consultar historial y detalle de pagos como cliente & MR-01, MR-07 & Historial y detalle de pagos del cliente \\ \midrule
		HU-27 & Consultar pagos recibidos como empresa & MR-01, MR-07 & Consulta de pagos por empresa \\ \midrule
		HU-28 & Cancelar pago simulado con motivo & MR-02, MR-07 & Cancelación permitida de pago simulado \\ \midrule
		HU-29 & Consultar indicadores de los cuatro perfiles & MR-01, MR-07 & Indicadores filtrados por perfil \\ \midrule
		HU-30 & Definir un recorrido con múltiples destinos & MR-05, MR-07 & Definición de recorrido con múltiples destinos \\ \midrule
		HU-31 & Calcular una ruta optimizada & MR-05, MR-07 & Cálculo de una ruta propuesta \\ \midrule
		HU-32 & Visualizar y guardar la ruta propuesta & MR-03, MR-07 & Visualización y persistencia de ruta planificada \\ \midrule
		HU-33 & Comparar los métodos de optimización & MR-05, MR-07 & Comparación experimental de métodos \\ \midrule
	\end{longtable}
\endgroup
\textit{Nota.} Elaboración propia a partir de las historias de usuario de TravelDLS. Los vínculos MR–HU son una propuesta de trazabilidad elaborada a partir de los resúmenes y descripciones disponibles.

\subsubsection{Matriz de criterios de aceptación}

Los criterios siguientes sintetizan las condiciones previstas para las historias de usuario de TravelDLS. Constituyen condiciones previstas para comprobar el comportamiento de las historias y no resultados de pruebas ya ejecutadas. En la implementación, los criterios deberán probarse con casos válidos, datos inválidos, permisos insuficientes y escenarios de error cuando corresponda.


\begingroup
	\small
	\setlength{\tabcolsep}{4.5pt}
	\renewcommand{\arraystretch}{1.45}
	\begin{longtable}{@{}>{\RaggedRight\arraybackslash}p{1.7cm}>{\RaggedRight\arraybackslash}p{3.9cm}>{\RaggedRight\arraybackslash}p{10.1cm}@{}}
		\caption{Matriz de criterios de aceptación}\label{tab:criterios_tactica}\\
		\toprule
		\textbf{Historia} & \textbf{Funcionalidad} & \textbf{Criterios de aceptación propuestos} \\
		\midrule
		\endfirsthead
		\caption[]{Matriz de criterios de aceptación (continuación)}\\
		\toprule
		\textbf{Historia} & \textbf{Funcionalidad} & \textbf{Criterios de aceptación propuestos} \\
		\midrule
		\endhead
		\endfoot
		\bottomrule
		\endlastfoot
		HU-01 & Cimientos de base de datos y backend & Preparar el backend y la persistencia de la aplicación. Configurar AdonisJS y PostgreSQL. Crear migraciones y modelos con sus relaciones. Verificar arranque, conexión y operaciones básicas. \\ \midrule
		HU-02 & Registrarse como cliente & Crear una cuenta con nombre, correo y contraseña válidos, siempre con rol cliente. Informar campos inválidos o correo duplicado sin crear cuentas repetidas. Acceder al área de cliente al completar el registro. \\ \midrule
		HU-03 & Iniciar, conservar y cerrar sesión & Acceder con credenciales válidas y mostrar error si son incorrectas. Conservar la sesión al recargar y solicitar acceso cuando expire. Cerrar sesión e impedir volver a consultar contenido privado. \\ \midrule
		HU-04 & Acceder a navegación y operaciones por rol & Mostrar las áreas y acciones correspondientes a cada rol. Impedir el acceso directo a áreas de otro perfil. Rechazar operaciones no autorizadas también desde la API. \\ \midrule
		HU-05 & Registrar usuarios con rol desde administración & Registrar nombre, correo, contraseña y un rol permitido. Informar datos inválidos o correo duplicado. Guardar el usuario con el rol elegido y confirmar el registro. \\ \midrule
		HU-06 & Administrar empresas & Consultar y buscar empresas en un listado paginado. Registrar y actualizar sus datos, usuario asociado y fotografía. Confirmar la baja antes de retirar una empresa del listado. \\ \midrule
		HU-07 & Consultar y dar de baja clientes desde administración & Consultar clientes y sus datos en un listado paginado. Solicitar confirmación antes de dar de baja a un cliente. Actualizar el listado tras la baja y mostrar errores si falla. \\ \midrule
		HU-08 & Completar y actualizar perfil de cliente & Consultar y completar los datos del perfil propio. Actualizar datos y fotografía, validando la información ingresada. Conservar los cambios después de recargar. \\ \midrule
		HU-09 & Actualizar perfil y tarifa de empresa & Consultar los datos de la empresa asociada a la sesión. Actualizar datos, fotografía y tarifa por kilogramo. Conservar los cambios e informar valores inválidos. \\ \midrule
		HU-10 & Completar y actualizar perfil de conductor & Consultar y completar el perfil del conductor. Actualizar datos y fotografía respetando las validaciones de identificación y licencia. Conservar los cambios y mostrar errores de validación. \\ \midrule
		HU-11 & Registrar, vincular y administrar conductores & Registrar conductores y vincularlos a una empresa según los permisos del usuario. Consultar y mantener los conductores mediante las acciones habilitadas para cada rol. Rechazar datos inválidos y operaciones sobre empresas no autorizadas. \\ \midrule
		HU-12 & Administrar categorías de vehículos & Consultar categorías y su estado en un listado paginado. Crear y editar categorías con datos válidos. Confirmar la baja y actualizar las opciones de los formularios de camiones. \\ \midrule
		HU-13 & Administrar camiones y asignarlos a conductores & Registrar y editar placa, chasis, categoría, empresa, conductor y estado. Impedir asignar un conductor que ya tenga otro camión activo. Consultar la flota y confirmar la baja de un camión. \\ \midrule
		HU-14 & Consultar camión asignado & Mostrar los datos del camión asignado al conductor autenticado. Indicar cuando no existe un camión asignado. Mostrar los errores de consulta sin presentar datos de otro conductor. \\ \midrule
		HU-15 & Reportar mantenimiento del camión & Reportar mantenimiento únicamente del camión asignado. Guardar el estado de mantenimiento y mostrarlo al recargar. Informar si no existe asignación o si falla la operación. \\ \midrule
		HU-16 & Solicitar transporte con empresa, carga y direcciones & Seleccionar empresa e indicar carga, cantidad, peso, embalaje y direcciones. Guardar el pedido y sus detalles, validando los datos obligatorios. Informar fallos o registros parciales sin duplicar la solicitud. \\ \midrule
		HU-17 & Consultar y gestionar pedidos del cliente & Consultar pedidos propios con estado, importe y detalle. Permitir cambios y bajas únicamente en los estados habilitados. Actualizar los datos después de cada operación y respetar filtros y paginación. \\ \midrule
		HU-18 & Consultar solicitudes recibidas por empresa & Consultar las solicitudes correspondientes a la empresa. Mostrar cliente, carga, estado e importe con filtros y paginación. Reflejar cambios de estado e informar cuando no existan solicitudes. \\ \midrule
		HU-19 & Establecer el precio de un pedido & Asignar un importe válido a un pedido de la empresa. Dejar el pedido pendiente de aprobación del cliente. Mostrar al cliente el precio actualizado e informar si la operación falla. \\ \midrule
		HU-20 & Responder a la cotización de transporte & Consultar el precio propuesto antes de decidir. Rechazar la cotización y reflejar el estado denegado. Confirmar la aceptación mediante el flujo de pago simulado, sin cambios de estado incompatibles. \\ \midrule
		HU-21 & Asignar conductor a un pedido & Seleccionar un conductor autorizado para el pedido. Rechazar la asignación al alcanzar cinco viajes activos según la regla vigente. Mostrar el pedido entre los viajes del conductor asignado. \\ \midrule
		HU-22 & Consultar, iniciar y finalizar viajes & Consultar únicamente los viajes asignados al conductor. Registrar inicio y entrega, actualizando la disponibilidad del conductor. Impedir iniciar pedidos anulados y detener el seguimiento al finalizar. \\ \midrule
		HU-23 & Compartir y consultar ubicación del viaje & Compartir la ubicación del conductor durante el viaje. Permitir que cliente y empresa participantes consulten la última posición. Informar falta de señal, permiso GPS o ubicación, sin exponer viajes ajenos. \\ \midrule
		HU-24 & Visualizar simulación del recorrido & Identificar el recorrido como una simulación, separada de la ubicación real. Permitir iniciar, pausar y reiniciar la animación con progreso y tiempo estimado. No modificar la ubicación real ni el estado del pedido. \\ \midrule
		HU-25 & Realizar pago simulado por tarjeta o transferencia & Simular el pago de un pedido propio pendiente de aprobación y con importe válido. Validar tarjeta o transferencia e impedir pagos duplicados. Guardar la referencia y aceptar el pedido, indicando que no existe un cobro real. \\ \midrule
		HU-26 & Consultar historial y detalle de pagos como cliente & Consultar pagos propios con estado, método, importe y referencia. Filtrar, paginar y abrir el detalle de un pago. Mostrar totales según el alcance consultado y gestionar listas vacías o errores. \\ \midrule
		HU-27 & Consultar pagos recibidos como empresa & Consultar únicamente los pagos asociados a la empresa. Filtrar, paginar y abrir el detalle de las operaciones. Mostrar importes y estados de los pagos simulados. \\ \midrule
		HU-28 & Cancelar pago simulado con motivo & Solicitar motivo válido y confirmación para cancelar un pago propio. Impedir cancelar pagos ya cancelados o de pedidos en tránsito o entregados. Guardar la cancelación y anular el pedido si estaba aceptado. \\ \midrule
		HU-29 & Consultar indicadores de los cuatro perfiles & Mostrar a cada perfil sus indicadores y accesos rápidos. Calcular los valores con datos reales y respetar el alcance de la consulta. Distinguir ausencia de actividad de errores de consulta y permitir vista móvil. \\ \midrule
		HU-30 & Definir un recorrido con múltiples destinos & Permitir agregar, editar y eliminar destinos antes de calcular la ruta. Validar coordenadas y datos obligatorios, indicando puntos inválidos o repetidos. Guardar el recorrido asociado a la empresa y limitar su acceso a usuarios autorizados. \\ \midrule
		HU-31 & Calcular una ruta optimizada & Visitar todos los destinos válidos respetando el origen y la condición de finalización definidos. Mostrar el orden recomendado y el costo total, con la unidad utilizada. Informar entradas incompletas, destinos inaccesibles o ausencia de solución sin mostrar resultados engañosos. \\ \midrule
		HU-32 & Visualizar y guardar la ruta propuesta & Representar el recorrido y los destinos numerados en el orden calculado. Guardar la propuesta únicamente tras la confirmación del usuario y permitir consultarla nuevamente. Distinguir la ruta planificada de la ubicación real y de la simulación del recorrido. \\ \midrule
		HU-33 & Comparar los métodos de optimización & Comparar Vecino Más Cercano, su mejora con 2-opt y OR-Tools con los mismos puntos, matriz y condiciones. Mostrar por método el recorrido, costo total y tiempo de cálculo. Conservar entradas, parámetros y resultados para repetir la evaluación, sin asumir un método ganador. \\ \midrule
	\end{longtable}
\endgroup
\textit{Nota.} HU-01 reúne las tareas iniciales de configuración y comprobación de la persistencia.

\subsubsection{Planificación propuesta de los sprints}

La Tabla \ref{tab:sprints} presenta una agrupación funcional inicial de las historias de usuario. Las fechas y la duración de los sprints deben validarse con el cronograma general antes de establecerse como planificación definitiva.


\begingroup
	\small
	\setlength{\tabcolsep}{4.5pt}
	\renewcommand{\arraystretch}{1.45}
	\begin{longtable}{@{}>{\RaggedRight\arraybackslash}p{1.9cm}>{\RaggedRight\arraybackslash}p{3.45cm}>{\RaggedRight\arraybackslash}p{2.6cm}>{\RaggedRight\arraybackslash}p{7.45cm}@{}}
		\caption{Planificación propuesta de los sprints}\label{tab:sprints}\\
		\toprule
		\textbf{Iteración} & \textbf{Objetivo} & \textbf{Historias} & \textbf{Entregable previsto} \\
		\midrule
		\endfirsthead
		\caption[]{Planificación propuesta de los sprints (continuación)}\\
		\toprule
		\textbf{Iteración} & \textbf{Objetivo} & \textbf{Historias} & \textbf{Entregable previsto} \\
		\midrule
		\endhead
		\endfoot
		\bottomrule
		\endlastfoot
		Sprint 1 & Fundamentos y acceso & HU-01 a HU-04 & Backend y PostgreSQL configurados; registro, autenticación, sesión y control de acceso básico comprobados. \\ \midrule
		Sprint 2 & Administración y perfiles & HU-05 a HU-10 & Gestión de usuarios, empresas y clientes; actualización de perfiles con validaciones. \\ \midrule
		Sprint 3 & Conductores y flota & HU-11 a HU-15 & Gestión de conductores, categorías y camiones; asignaciones, disponibilidad y reporte de mantenimiento. \\ \midrule
		Sprint 4 & Pedidos y cotizaciones & HU-16 a HU-20 & Solicitud, consulta, cotización y respuesta del cliente; reglas de estados y visualización del importe. \\ \midrule
		Sprint 5 & Asignación, viajes y seguimiento & HU-21 a HU-24 & Asignación de conductor, inicio y finalización de viajes, ubicación real y simulación separada. \\ \midrule
		Sprint 6 & Pagos simulados e indicadores & HU-25 a HU-29 & Pago simulado, historial y consultas por perfil, cancelación permitida e indicadores con permisos. \\ \midrule
		Sprint 7 & Planificación y comparación de rutas & HU-30 a HU-33 & Definición de destinos, cálculo, mapa, persistencia y comparación de métodos con entradas equivalentes. \\ \midrule
	\end{longtable}
\endgroup
\textit{Nota.} Elaboración propia.


\subsubsection{Construcción del prototipo}

La construcción de TravelDLS se organiza mediante las historias de usuario y los sprints propuestos. En el backend se desarrollaron rutas y servicios para autenticación, administración de usuarios y recursos de transporte, gestión de órdenes, cotización, asignación de conductores, pagos simulados, consulta de ubicación y confirmación de entregas. En el frontend se incorporaron interfaces para administración, clientes, empresas transportistas y conductores, con pantallas para consultar recursos y pedidos, solicitar transporte, responder a cotizaciones y consultar información del viaje. La disponibilidad de estos componentes no supone que todos los recorridos entre interfaz y servicios se hayan verificado.

Para el seguimiento se incorporó el envío de la ubicación reportada por el dispositivo del conductor y su consulta por participantes autorizados. La visualización utiliza OpenStreetMap y Leaflet, y durante el desarrollo el trazado de rutas puede consultar un servicio de prueba basado en OSRM. El backend conserva la última posición asociada al pedido, pero aún no construye un historial de puntos. La incorporación de Mapbox y del historial requiere adaptar las vistas y definir el almacenamiento de ubicaciones sucesivas. También se incorporó código para gestionar con Cloudinary fotografías de perfiles, empresas y entregas; las pruebas unitarias validan los archivos antes de su carga, pero falta comprobar el envío al servicio externo.

En la contratación se desarrollaron el registro de precios de las cargas, la respuesta del cliente a la cotización y un pago simulado por tarjeta o transferencia. Las pruebas unitarias comprueban el cálculo de montos y algunas validaciones del medio de pago; una prueba funcional comprueba el registro de una transferencia simulada y el rechazo de un segundo pago para la misma orden. Sigue pendiente verificar el recargo del 8\,\% y la secuencia completa de cotización, aceptación y pago. Se dispone de documentación técnica del frontend, los mapas y la confirmación de entregas, así como de archivos de prueba para contratos del frontend; los manuales de usuario todavía no se han documentado.

\subsection{Cronograma de actividades}

El cronograma de trabajo presenta las actividades desarrolladas durante
el proceso de documentación y desarrollo del software TravelDLS,
indicando el período establecido para cada una y su estado de avance.

\begin{center}
	\begin{ganttchart}[
		hgrid,
		vgrid,
		x unit=1.35cm,
		y unit title=0.60cm,
		y unit chart=0.42cm,
		title height=1,
		bar height=0.24,
		group height=0.28,
		title label font=\scriptsize,
		bar label font=\scriptsize,
		group label font=\bfseries\scriptsize,
		title/.style={
			fill=white,
			draw=grislinea
		},
		group/.style={
			fill=white,
			draw=none
		},
		bar/.style={
			draw=none
		},
		bar label node/.append style={
			align=right,
			text width=6.1cm,
			anchor=east,
			xshift=-0.2cm
		},
		group label node/.append style={
			align=right,
			text width=6.1cm,
			anchor=east,
			xshift=-0.2cm
		}
		]{1}{7}
		
		% ENCABEZADO DE SEMANAS
		\gantttitle{\shortstack{\scriptsize Semana 1\\\tiny 7--11 sep}}{1}
		\gantttitle{\shortstack{\scriptsize Semana 2\\\tiny 14--18 sep}}{1}
		\gantttitle{\shortstack{\scriptsize Semana 3\\\tiny 21--25 sep}}{1}
		\gantttitle{\shortstack{\scriptsize Semana 4\\\tiny 28 sep--2 oct}}{1}
		\gantttitle{\shortstack{\scriptsize Semana 5\\\tiny 5--9 oct}}{1}
		\gantttitle{\shortstack{\scriptsize Semana 6\\\tiny 12--16 oct}}{1}
		\gantttitle{\shortstack{\scriptsize Semana 7\\\tiny 19--23 oct}}{1}\\
		
		% DOCUMENTACIÓN
		\ganttgroup{DOCUMENTACIÓN}{1}{7}\\
		
		\ganttbar[bar/.style={fill=completado,draw=none}]
		{Revisión bibliográfica}{1}{1}\\
		
		\ganttbar[bar/.style={fill=completado,draw=none}]
		{Objetivos y tema}{1}{1}\\
		
		\ganttbar[bar/.style={fill=completado,draw=none}]
		{Planteamiento, alcance y limitaciones}{2}{2}\\
		
		\ganttbar[bar/.style={fill=completado,draw=none}]
		{Marco teórico}{2}{2}\\
		
		\ganttbar[bar/.style={fill=completado,draw=none}]
		{Factibilidad del proyecto}{3}{3}\\
		
		\ganttbar[bar/.style={fill=completado,draw=none}]
		{Metodología}{4}{6}\\
		
		\ganttbar[bar/.style={fill=proceso,draw=none}]
		{Plan de ejecución}{6}{6}\\
		
		\ganttbar[bar/.style={fill=proceso,draw=none}]
		{Resultados del plan de ejecución}{7}{7}\\
		
		% BACKEND
		\ganttgroup{BACKEND}{1}{6}\\
		
		\ganttbar[bar/.style={fill=completado,draw=none}]
		{Diagrama entidad-relación (DER)}{1}{1}\\
		
		\ganttbar[bar/.style={fill=completado,draw=none}]
		{Usuarios, roles y compañía -- CRUD}{2}{2}\\
		
		\ganttbar[bar/.style={fill=completado,draw=none}]
		{Clientes, conductores y camiones -- CRUD}{3}{3}\\
		
		\ganttbar[bar/.style={fill=completado,draw=none}]
		{Categorías de camiones y órdenes -- CRUD}{4}{4}\\
		
		\ganttbar[bar/.style={fill=completado,draw=none}]
		{Autenticación JWT y pruebas de integración}{4}{4}\\
		
		\ganttbar[bar/.style={fill=completado,draw=none}]
		{Migraciones, modelos y validaciones}{5}{5}\\
		
		\ganttbar[bar/.style={fill=proceso,draw=none}]
		{Finalización de CRUD y pruebas del backend}{6}{6}\\
		
		% FRONTEND
		\ganttgroup{FRONTEND}{7}{7}\\
		
		\ganttbar[bar/.style={fill=proceso,draw=none}]
		{Configuración de Next.js y autenticación JWT}{7}{7}
		
	\end{ganttchart}
	
	\vspace{0.3cm}
	
	\begin{flushleft}
		\scriptsize
		\textbf{Estado:}
		\quad
		\textcolor{completado}{\rule{0.45cm}{0.16cm}}
		\; Completado
		\qquad
		\textcolor{proceso}{\rule{0.45cm}{0.16cm}}
		\; En proceso
	\end{flushleft}
\end{center}

\subsection{Plan de ejecución}

El plan de ejecución organiza las acciones de análisis, diseño y evaluación previstas para los tres objetivos específicos de TravelDLS, junto con sus instrumentos de verificación, protagonistas y resultados esperados. La construcción de componentes y las integraciones técnicas se describen en la metodología de desarrollo. Para el primer objetivo se contempla una revisión documental con cadenas de búsqueda en Google Académico, tomando como referencia el mapeo sistemático de \textcite{petersen2008systematic}. La identificación de requisitos se fundamenta en la documentación y los modelos del proyecto; no corresponde a un levantamiento formal con usuarios.

	\begingroup
	\scriptsize
	\setlength{\tabcolsep}{2pt}
	\renewcommand{\arraystretch}{1.15}
	\begin{longtable}{@{}>{\RaggedRight\arraybackslash}p{1.5cm}>{\RaggedRight\arraybackslash}p{2.3cm}>{\RaggedRight\arraybackslash}p{3.75cm}>{\RaggedRight\arraybackslash}p{2.9cm}>{\RaggedRight\arraybackslash}p{2.1cm}>{\RaggedRight\arraybackslash}p{2.85cm}@{}}
	\caption{Plan de ejecución de TravelDLS}\label{tab:plan_ejecucion}\\
	\toprule
	\textbf{Objetivos} & \textbf{Acciones} & \textbf{Descripción de las acciones} & \textbf{Instrumento de verificación} & \textbf{Protagonistas} & \textbf{Resultados esperados} \\
	\midrule
	\endfirsthead
	\caption[]{Plan de ejecución de TravelDLS (continuación)}\\
	\toprule
	\textbf{Objetivos} & \textbf{Acciones} & \textbf{Descripción de las acciones} & \textbf{Instrumento de verificación} & \textbf{Protagonistas} & \textbf{Resultados esperados} \\
	\midrule
	\endhead
	\endfoot
	\bottomrule
	\endlastfoot
	\hypertarget{plan-oe1}{OE1} & Revisión documental y diagnóstico & Seleccionar antecedentes, examinar los procesos de transporte de carga pesada y delimitar necesidades de gestión, seguimiento y planificación de rutas. & Cadenas de búsqueda aplicadas en Google Académico, tomando como referencia el procedimiento de mapeo sistemático (Systematic Mapping) de Petersen et al. (2008). & Equipo investigador. & Problemática fundamentada y necesidades iniciales identificadas. \\ \midrule
	OE1 & Identificación de requisitos del sistema & Analizar los objetivos, el alcance y los modelos de TravelDLS para identificar las funcionalidades y necesidades técnicas de la plataforma. No se realizó levantamiento directo con usuarios. & Documento del proyecto, diagramas UML y entidad-relación. & Equipo investigador. & Se identificaron preliminarmente las funcionalidades y componentes de TravelDLS mediante análisis documental, sin validación directa con usuarios. \\ \midrule
	OE1 & Análisis de métodos de rutas & Definir un escenario hipotético con múltiples destinos, calcular recorridos con Vecino Más Cercano y 2-opt, y estudiar la formulación de Google OR-Tools para una comparación posterior bajo las mismas condiciones. & Matriz de distancias hipotética, cálculos reproducibles y criterios de comparación. & Equipo investigador. & Fundamentos y datos de referencia para evaluar los métodos cuando OR-Tools se integre; sin atribuirle resultados de ejecución. \\ \midrule
	OE2 & Definición de arquitectura & Establecer la organización del frontend con Next.js, el backend con AdonisJS bajo MVC y capa de servicios, la base de datos PostgreSQL y los mecanismos de acceso por empresa. & Documento de arquitectura y esquema de componentes. & Equipo de desarrollo. & Arquitectura y responsabilidades técnicas de los componentes definidas. \\ \midrule
	\hypertarget{plan-oe2}{OE2} & Modelado de datos y clases & Elaborar y revisar los diagramas de clases y entidad-relación, considerando usuarios, roles, empresas, vehículos, conductores, clientes, órdenes y pagos. & Diagramas UML y entidad-relación revisados. & Equipo de desarrollo. & Modelos coherentes con las entidades y relaciones previstas en TravelDLS. \\ \midrule
	OE2 & Modelado de procesos & Revisar los diagramas de casos de uso, estados y secuencia, unificando las reglas sobre cotización, pago, asignación del conductor y seguimiento de las órdenes. & Diagramas de casos de uso, estados y secuencia; registro de revisión. & Equipo de desarrollo y asesoría académica. & Flujos de negocio y responsabilidades consistentes entre los modelos. \\ \midrule
	\hypertarget{plan-oe3}{OE3} & Diseño conceptual del módulo de optimización & Especificar los destinos, costos, restricciones, entradas y salidas del modelo de rutas con Google OR-Tools para orientar su futura integración, sin presentar el módulo como implementado. & Formulación del problema, matriz de distancias y criterios de evaluación. & Equipo investigador y de desarrollo. & Diseño documentado y compromiso de integrar y evaluar el módulo en una etapa posterior. \\ \midrule
	OE3 & Evaluación preliminar del prototipo & Definir escenarios y criterios de prueba, ejecutar los casos disponibles y registrar los resultados observados y los aspectos aún no verificados. & Casos de prueba, registro de ejecución y matriz de resultados. & Equipo de desarrollo y participantes de las pruebas. & Evidencia delimitada de los comportamientos comprobados y de las verificaciones pendientes. \\
	\end{longtable}
	\endgroup

\clearpage
\subsection{Resultados del plan de ejecución}
\label{sec:resultados-plan-ejecucion}

Los resultados del análisis, diseño y evaluación preliminar de TravelDLS se presentan de acuerdo con las acciones establecidas en el plan de ejecución. En relación con el \hyperlink{plan-oe1}{OE1}, se realizó la revisión documental y la identificación preliminar de los requisitos del sistema, lo que permitió fundamentar la problemática e identificar las necesidades iniciales de la plataforma. Para orientar la búsqueda de información, se definieron términos en español e inglés relacionados con el seguimiento y rastreo de vehículos y carga, el transporte de mercancías, la logística, los envíos y la optimización de rutas. Asimismo, se analizaron los objetivos, el alcance y los modelos del proyecto para establecer las funcionalidades y los componentes principales de TravelDLS, sin realizar un levantamiento directo con usuarios.

Las cadenas de búsqueda utilizadas para la revisión documental combinan palabras clave y operadores booleanos, como \texttt{AND} y \texttt{OR}, con el propósito de orientar la consulta de información relacionada con el transporte de carga, el seguimiento de vehículos, la optimización de rutas y las herramientas tecnológicas consideradas para el desarrollo de TravelDLS.

\textbf{Cadenas de búsqueda utilizadas}

\textbf{1. Seguimiento y transporte de carga}

\begin{quote}
	\begin{lstlisting}[style=cadenas]
		("tracking" OR "seguimiento" OR "rastreo" OR "monitoreo") AND ("carga" OR "transporte de carga" OR "freight" OR "freight transportation" OR "logistics")
	\end{lstlisting}
\end{quote}

\begin{quote}
	\begin{lstlisting}[style=cadenas]
		("freight transportation" OR "heavy freight" OR "cargo transportation" OR "logistics") AND ("tracking" OR "real-time tracking" OR "vehicle tracking" OR "cargo tracking") AND ("resource management" OR "fleet management" OR "transport management" OR "logistics management")
	\end{lstlisting}
\end{quote}

Las cadenas de búsqueda de este apartado se utilizaron para localizar investigaciones relacionadas con el seguimiento de cargas, la localización de vehículos y la gestión del transporte. La información obtenida sirvió principalmente para elaborar los antecedentes del proyecto, identificando sistemas e investigaciones que abordan la administración de fletes, el seguimiento en tiempo real y la coordinación de las operaciones logísticas. También contribuyó al marco teórico, especialmente en la definición de conceptos relacionados con la gestión logística, el seguimiento de cargas y la localización de vehículos.

\textbf{2. Optimización de rutas y algoritmos}

\begin{quote}
	\begin{lstlisting}[style=cadenas]
		("traveling salesman problem" OR TSP OR "vehicle routing problem" OR VRP) AND ("nearest neighbor" OR "nearest neighbour" OR "2-opt" OR "OR-Tools") AND ("logistics" OR "freight transportation" OR "delivery routing")
	\end{lstlisting}
\end{quote}


\begin{quote}
	\begin{lstlisting}[style=cadenas]
		("2-opt" OR "two-opt" OR "2 opt") AND ("traveling salesman problem" OR "vehicle routing problem" OR "route optimization") AND ("logistics" OR "transportation" OR "freight transport")
	\end{lstlisting}
\end{quote}
Estas cadenas se utilizaron para buscar investigaciones sobre la planificación y optimización de rutas, con énfasis en el Problema del Agente Viajero (TSP), el algoritmo Vecino Más Cercano y el método 2-opt. Los resultados se emplearon en los antecedentes para identificar trabajos que han aplicado métodos de optimización en el transporte y la distribución de mercancías. Asimismo, sirvieron para desarrollar las bases teóricas del proyecto, explicar los métodos considerados y fundamentar la comparación de alternativas para la planificación de recorridos en TravelDLS.

\textbf{3. Google OR-Tools y matrices de distancias}

\begin{quote}
	\begin{lstlisting}[style=cadenas]
		("Google OR-Tools" OR "OR-Tools" OR "Google Operations Research Tools") AND ("vehicle routing problem" OR "traveling salesman problem" OR "route optimization") AND ("logistics" OR "transportation" OR "delivery")
	\end{lstlisting}
\end{quote}

\begin{quote}
	\begin{lstlisting}[style=cadenas]
		("distance matrix" OR "travel time matrix" OR "road network distances") AND ("route optimization" OR "vehicle routing problem" OR "traveling salesman problem") AND ("logistics" OR "freight transportation" OR "delivery routing")
	\end{lstlisting}
\end{quote}

Las cadenas de este apartado se utilizaron para investigar el funcionamiento y la aplicación de Google OR-Tools en problemas de optimización de rutas, así como el uso de matrices de distancias y tiempos de recorrido. La información recopilada contribuyó principalmente a las bases teóricas y a la fundamentación técnica del módulo de optimización de TravelDLS. Además, permitió identificar elementos necesarios para comparar diferentes recorridos y comprender cómo estas herramientas pueden apoyar la planificación de viajes con múltiples destinos.

\textbf{4. Geolocalización y seguimiento mediante Mapbox}

\begin{quote}
	\begin{lstlisting}[style=cadenas]
		("Mapbox" OR "Mapbox API") AND ("truck tracking" OR "vehicle tracking" OR "fleet tracking" OR "GPS tracking") AND ("logistics" OR "freight transportation" OR "cargo transportation" OR "delivery")
	\end{lstlisting}
\end{quote}

\begin{quote}
	\begin{lstlisting}[style=cadenas]
		("Mapbox" OR "Mapbox API" OR "Mapbox SDK") AND ("real-time tracking" OR "vehicle tracking" OR "GPS tracking" OR geolocation) AND (logistics OR transportation OR "fleet management" OR delivery OR travel)
	\end{lstlisting}
\end{quote}

La consulta de documentación sobre Mapbox, sus API y SDK permitió conocer las herramientas disponibles para trabajar con mapas, geolocalización y seguimiento de vehículos. La información obtenida se consideró en la selección de tecnologías y en el análisis de la viabilidad técnica del proyecto, como parte de la evaluación de las herramientas que podrían apoyar las funcionalidades de ubicación y seguimiento de recorridos.

\textbf{5. Almacenamiento de imágenes mediante Cloudinary}

\begin{quote}
	\begin{lstlisting}[style=cadenas]
		("Cloudinary" OR "Cloudinary API") AND ("web application" OR "mobile application" OR platform OR system) AND (logistics OR transportation OR travel OR delivery OR tracking OR fleet) AND ("image storage" OR "image management" OR "media storage" OR "cloud storage")
	\end{lstlisting}
\end{quote}

Las búsquedas sobre Cloudinary y el almacenamiento de imágenes y recursos multimedia en aplicaciones web permitieron consultar información sobre servicios para gestionar este tipo de archivos. Los resultados se consideraron en la selección de tecnologías y en el análisis de la viabilidad técnica de TravelDLS, con el propósito de evaluar una alternativa para el almacenamiento y la gestión de recursos gráficos asociados al sistema.

\paragraph{Ejercicios preliminares de optimización de rutas}
Para la acción «Análisis de métodos de rutas» del OE1 se prepararon dos ejercicios \emph{hipotéticos y reproducibles}: un camino entre una recogida y una entrega, como ilustración complementaria, y un recorrido con cuatro destinos, relacionado directamente con el escenario del plan. Estos ejercicios permiten documentar los cálculos de Vecino Más Cercano y 2-opt y establecer una matriz común para la comparación futura con Google OR-Tools. Las distancias siguientes pertenecen únicamente a los datos ficticios de cada ejercicio; no son distancias reales ni resultados obtenidos por OR-Tools.

\textbf{Ejercicio complementario: trayecto directo A--B.} Se consideraron cuatro nodos: A (recogida), B (entrega) y dos intersecciones, X e Y. Las conexiones disponibles y sus distancias hipotéticas son A--X (4 km), X--B (6 km), A--Y (3 km), Y--B (10 km) y A--B (12 km). Se supone que pueden recorrerse en los sentidos indicados. La Tabla~\ref{tab:caminos-hipoteticos} compara los tres caminos posibles de esta red simplificada.

\begin{table}[H]
\centering
\caption{Caminos del ejercicio directo A--B}
\label{tab:caminos-hipoteticos}
\begin{tabular}{lcc}
\toprule
Camino & Cálculo & Distancia \\
\midrule
A--X--B & $4+6$ & 10 km \\
A--B & $12$ & 12 km \\
A--Y--B & $3+10$ & 13 km \\
\bottomrule
\end{tabular}
\end{table}

El menor costo es 10 km para A--X--B. Este valor se obtuvo por enumeración de los caminos del ejemplo; queda pendiente comprobar la salida de \texttt{SimpleMinCostFlow} con la misma red. En este caso no hay un orden de entregas que reorganizar, por lo que Vecino Más Cercano y 2-opt no se aplican a la comparación.

\textbf{Ejercicio principal: una recogida y cuatro entregas.} Se fijó A como punto de recogida inicial y B, C, D y E como destinos obligatorios. Para este ejercicio se supone un solo vehículo, distancias simétricas y retorno a A. La Tabla~\ref{tab:matriz-rutas-hipotetica} muestra la matriz ficticia, expresada en kilómetros.

\begin{table}[H]
\centering
\caption{Matriz de distancias hipotéticas para cinco ubicaciones}
\label{tab:matriz-rutas-hipotetica}
\begin{tabular}{crrrrr}
\toprule
Nodo & A & B & C & D & E \\
\midrule
A & 0 & 7 & 8 & 9 & 13 \\
B & 7 & 0 & 1 & 7 & 7 \\
C & 8 & 1 & 0 & 7 & 7 \\
D & 9 & 7 & 7 & 0 & 14 \\
E & 13 & 7 & 7 & 14 & 0 \\
\bottomrule
\end{tabular}
\end{table}

Al aplicar la técnica de Vecino Más Cercano \parencite{contreras2021aplicacion}, desde A se elige B (7 km) y luego C (1 km). Ante el empate de 7 km entre D y E desde C, se adopta el desempate alfabético y se elige D. Con los datos ficticios de la Tabla~\ref{tab:matriz-rutas-hipotetica}, el recorrido resultante A--B--C--D--E--A mide $7+1+7+14+13=42$ km.

Para mejorar ese recorrido se aplica la heurística 2-opt, que intercambia dos aristas e invierte el segmento intermedio \parencite{dacosta2021twoopt,uddin2023twoopt}. En este ejemplo se reemplazan A--B y D--E, cuyo costo conjunto es 21 km, por A--D y B--E, cuyo costo es 16 km. Según la misma matriz ficticia, la inversión produce A--D--C--B--E--A, con una distancia de $9+7+1+7+13=37$ km. La reducción respecto de la ruta inicial es $(42-37)/42\times100=11{,}90$\,\%.

La enumeración de las $4!=24$ secuencias de los destinos, manteniendo fijo A al inicio y al final, confirma que 37 km es el mínimo de esta matriz; la ruta mostrada es una de las que alcanzan ese valor. Esto no demuestra que 2-opt encuentre siempre un mínimo global. La Tabla~\ref{tab:comparacion-rutas-preliminar} separa los cálculos ya comprobados de la evaluación aún pendiente.

\begin{table}[H]
\centering
\caption{Estado de la comparación del ejercicio con múltiples destinos}
\label{tab:comparacion-rutas-preliminar}
\begin{tabular}{@{}p{0.52\textwidth}p{0.18\textwidth}p{0.2\textwidth}@{}}
\toprule
Procedimiento & Distancia & Estado \\
\midrule
Vecino Más Cercano & 42 km & Cálculo manual \\
Vecino Más Cercano y 2-opt & 37 km & Cálculo manual \\
Enumeración de 24 recorridos & 37 km & Mínimo verificado \\
Google OR-Tools (\texttt{RoutingModel}) & Pendiente & Sin implementación \\
\bottomrule
\end{tabular}
\end{table}

Se asume el compromiso de integrar el módulo de optimización con Google OR-Tools en una etapa posterior del desarrollo de TravelDLS, conforme al objetivo específico del proyecto. Para completar la comparación se ejecutará \texttt{RoutingModel} sobre la misma matriz y condiciones, se registrarán el recorrido obtenido y su distancia, y se medirán los tiempos de cálculo bajo un entorno común. La demostración directa podrá comprobarse por separado con \texttt{SimpleMinCostFlow}. Hasta entonces, este resultado documenta una preparación matemática y criterios para la evaluación futura, sin atribuir al prototipo rutas optimizadas ni dar por cumplida la implementación del módulo.

En relación con el \hyperlink{plan-oe2}{OE2}, se definió una arquitectura que separa la interfaz desarrollada con Next.js, los controladores y servicios del backend construido con AdonisJS y la persistencia en PostgreSQL. Se elaboraron los modelos de clases y entidad-relación para representar usuarios, roles, empresas, vehículos, conductores, clientes, órdenes, cargas y pagos, así como sus asociaciones. Estos modelos permiten describir la organización prevista de los datos y las responsabilidades de los componentes. Aún queda pendiente documentar en un esquema de componentes la interacción entre las tres capas y comprobar de manera sistemática la correspondencia de los modelos con las estructuras implementadas y las reglas de acceso por empresa. Por ello, se presenta un diseño elaborado, sin afirmar que su revisión integral se haya completado.

Como evidencia de las acciones de «Modelado de datos y clases» y «Modelado de procesos» previstas para el OE2 en el plan de ejecución, se presentan a continuación los diagramas elaborados. Primero se muestran los modelos de la estructura y los datos del sistema; después, los modelos de sus procesos. Estas figuras documentan el avance del diseño, cuya correspondencia con la implementación todavía requiere una revisión sistemática.

\newcommand{\reservarEspacioFigura}[1]{%
	\par
	\ifdim\dimexpr\pagegoal-\pagetotal\relax<#1\relax
		\newpage
	\fi
}

\reservarEspacioFigura{0.62\textheight}
\begin{samepage}
En la Figura~\ref{fig:diagrama-clases} se presentan las clases principales de TravelDLS, sus atributos y las relaciones entre usuarios, empresas, clientes, conductores y órdenes. También se muestran los servicios encargados de la autenticación, la gestión de órdenes y los pagos simulados.

\begin{figure}[H]
	\centering
	\includegraphics[width=\textwidth,height=0.82\textheight,keepaspectratio]{diagrama-clases.png}
	\caption{Diagrama de clases de TravelDLS: modelado de datos y clases del OE2}
	\label{fig:diagrama-clases}
\end{figure}
\end{samepage}

\medskip
\reservarEspacioFigura{0.78\textheight}
\begin{samepage}
En la Figura~\ref{fig:diagrama-er} se presenta el diagrama entidad-relación de TravelDLS, que muestra las tablas de la base de datos y cómo se relacionan entre sí. Este diagrama permite comprender cómo se almacena la información del sistema y sirve como guía para organizar los datos durante el desarrollo.

\begin{figure}[H]
	\centering
	\includegraphics[width=\textwidth,height=0.82\textheight,keepaspectratio]{diagrama-er.png}
	\caption{Diagrama entidad-relación de TravelDLS: modelado de datos y clases del OE2}
	\label{fig:diagrama-er}
\end{figure}
\end{samepage}

Los diagramas de casos de uso, estados y secuencia corresponden a la acción «Modelado de procesos» del OE2. El primero presenta las funciones por actor; el segundo, los estados de una orden; y el tercero, las interacciones del flujo principal de transporte. En comparación con el resultado esperado de obtener flujos y responsabilidades consistentes entre los modelos, también se considera un cumplimiento parcial: los procesos previstos están representados, pero no se incorpora un registro de revisión comparativa que confirme la coherencia de todas sus reglas.

\medskip
\reservarEspacioFigura{0.84\textheight}
\begin{samepage}
En la Figura~\ref{fig:diagrama-casos-uso} se describen las funciones previstas para el cliente, la empresa transportista y el conductor, como solicitar transporte, gestionar recursos y reportar la ubicación del viaje. El actor «Usuario» agrupa las funciones compartidas de acceso y administración de la cuenta.

\begin{figure}[H]
	\centering
	\includegraphics[width=0.88\textwidth,height=0.82\textheight,keepaspectratio]{diagrama-casos-uso.png}
	\caption{Diagrama de casos de uso de TravelDLS: modelado de procesos del OE2}
	\label{fig:diagrama-casos-uso}
\end{figure}

\smallskip
La Figura~\ref{fig:diagrama-estados} representa los estados de una orden desde su creación hasta la entrega o anulación, incluidos el pago, la asignación del conductor y las transiciones ante fallos.

\begin{figure}[H]
	\centering
	\includegraphics[width=\textwidth,height=0.82\textheight,keepaspectratio]{diagrama-estado-actualizado.png}
	\caption{Diagrama de estados de una orden de transporte: modelado de procesos del OE2}
	\label{fig:diagrama-estados}
\end{figure}
\end{samepage}

\medskip
\reservarEspacioFigura{0.86\textheight}
\begin{samepage}
En la Figura~\ref{fig:diagrama-secuencia} se muestra el intercambio de solicitudes y respuestas entre el cliente, la empresa, el conductor, el frontend, el backend y la base de datos. Las interacciones se organizan en creación del pedido, aprobación del precio, pago simulado, asignación del conductor y seguimiento hasta la entrega.

\begin{figure}[H]
	\centering
	\includegraphics[width=\textwidth,height=0.75\textheight,keepaspectratio]{diagrama-secuencia.png}
	\caption{Diagrama de secuencia del proceso principal de una orden: modelado de procesos del OE2}
	\label{fig:diagrama-secuencia}
\end{figure}
\end{samepage}

\clearpage
En relación con el \hyperlink{plan-oe3}{OE3}, se formuló un diseño conceptual del módulo de optimización para recorridos con múltiples destinos. Se definieron un origen, una matriz de costos y la condición de visitar todos los destinos, así como la necesidad de registrar el orden de visita, la distancia total y el tiempo de cálculo cuando se ejecute la comparación. Este resultado corresponde a la especificación del modelo; Google OR-Tools aún no se integró al prototipo. Se mantiene el compromiso de implementar el módulo y comprobar sus recorridos en una etapa posterior, por lo que el objetivo de implementación no se presenta como cumplido.

Para la acción «Evaluación preliminar del prototipo» se incorporaron pruebas unitarias de reglas de entrega, cifrado, validación de imágenes y operaciones de pago, junto con pruebas funcionales de autenticación, recursos de transporte, órdenes, permisos y respuestas de la API. En una ejecución realizada el 9 de noviembre de 2026 con \texttt{pnpm test} sobre una base de datos aislada para pruebas, se registraron 14 casos unitarios y 33 funcionales aprobados, sin fallos. La Tabla~\ref{tab:pruebas-backend-oe3} resume el resultado observado; el registro completo de la ejecución se conserva como evidencia técnica.

\begin{table}[H]
\centering
\caption{Resultado de la evaluación preliminar del backend para el OE3}
\label{tab:pruebas-backend-oe3}
\begin{tabular}{@{}lrrr@{}}
\toprule
Suite & Aprobadas & Fallidas & Total \\
\midrule
Unitaria & 14 & 0 & 14 \\
Funcional & 33 & 0 & 33 \\
\midrule
Total & 47 & 0 & 47 \\
\bottomrule
\end{tabular}
\end{table}

La aprobación de los casos ejecutados documenta únicamente el comportamiento verificado en esos escenarios; no equivale a una evaluación integral del prototipo. Aún falta consolidar una matriz que relacione cada caso funcional con su resultado, fecha e incidencia, y ampliar las pruebas de flujos completos, navegación e integraciones externas. La evaluación del módulo de optimización también permanece pendiente hasta que se implemente Google OR-Tools.

\subsection{Compromisos para la siguiente etapa}
\label{sec:compromisos-siguiente-etapa}

Los ejercicios preliminares de rutas y la consulta actual de ubicaciones permiten delimitar el trabajo pendiente, pero no acreditan la integración de Google OR-Tools ni de Mapbox. Para completar estas capacidades se asumen los siguientes compromisos:

\begin{itemize}
	\item Se completará el módulo de optimización con Google OR-Tools para recorridos con múltiples destinos. La implementación utilizará la matriz de distancias y las condiciones definidas en los ejercicios preliminares. Después, se comparará el recorrido obtenido con Vecino Más Cercano y 2-opt bajo los mismos datos y se registrarán el orden de visita, la distancia total y el tiempo de cálculo de cada método. Esta comparación permitirá evaluar el módulo con evidencia reproducible, sin anticipar resultados que aún no se han comprobado.
	\item Se integrará Mapbox en la visualización y consulta de ubicaciones del seguimiento de viajes. Actualmente, esas vistas utilizan OpenStreetMap y Leaflet; por ello, será necesario adaptar el mapa para representar con Mapbox la posición reportada por el conductor y asociada a cada orden. Una vez realizada la integración, se comprobará que los participantes autorizados puedan consultar la ubicación correspondiente y que esta se muestre correctamente en la interfaz. Hasta completar esa verificación, Mapbox permanecerá como componente pendiente.
\end{itemize}
